\documentclass[12pt]{article}
% Спільна преамбула для звітів (XeLaTeX)
\usepackage{fontspec}
\usepackage{polyglossia}
\setmainlanguage{ukrainian}
\setotherlanguage{english}
\setmainfont{DejaVuSerif}[Extension=.ttf,UprightFont=*,BoldFont=*-Bold,ItalicFont=*-Italic,BoldItalicFont=*-BoldItalic]
\setsansfont{DejaVuSans}[Extension=.ttf,UprightFont=*,BoldFont=*-Bold,ItalicFont=*-Oblique,BoldItalicFont=*-BoldOblique]
\setmonofont{DejaVuSansMono}[Extension=.ttf,UprightFont=*,BoldFont=*-Bold,ItalicFont=*-Oblique,BoldItalicFont=*-BoldOblique,Scale=0.9]
\usepackage[a4paper,margin=2.2cm]{geometry}
\usepackage{booktabs}
\usepackage{tabularx}
\usepackage{array}
\usepackage{enumitem}
\usepackage{xcolor}
\usepackage{tcolorbox}
\usepackage{graphicx}
\usepackage{setspace}
\usepackage[hidelinks]{hyperref}
\usepackage{titlesec}
\usepackage{fancyhdr}
\usepackage{multirow}
\usepackage{colortbl}
\setlength{\parskip}{4pt}
\setlength{\parindent}{0pt}
\onehalfspacing

\definecolor{accent}{HTML}{1F4E79}
\definecolor{sgreen}{HTML}{E2F0D9}
\definecolor{wred}{HTML}{FBE5D6}
\definecolor{oblue}{HTML}{DEEBF7}
\definecolor{tyellow}{HTML}{FFF2CC}
\definecolor{lightgray}{HTML}{F2F2F2}

\titleformat{\section}{\large\bfseries\color{accent}}{\thesection}{0.6em}{}
\titleformat{\subsection}{\normalsize\bfseries}{\thesubsection}{0.6em}{}

\pagestyle{fancy}
\fancyhf{}
\lhead{\small ІТ, Модуль 2 — командний проєкт}
\rhead{\small Петров Б., Мекшаков А.}
\cfoot{\thepage}
\renewcommand{\headrulewidth}{0.3pt}

\newcolumntype{L}[1]{>{\raggedright\arraybackslash}p{#1}}
\newcolumntype{Y}{>{\raggedright\arraybackslash}X}

% Титульна сторінка
\newcommand{\maketitlepage}[3]{%
\begin{titlepage}
\centering
{\large Київський національний університет імені Тараса Шевченка\\[2pt]
Факультет комп'ютерних наук та кібернетики\par}
\vspace{3cm}
{\large Дисципліна: \textbf{Інформаційні технології}\par}
\vspace{0.5cm}
{\large Модуль 2 — робота в команді\par}
\vspace{2cm}
{\LARGE\bfseries #1\par}
\vspace{0.6cm}
{\large #2\par}
\vfill
\begin{flushright}
\begin{tabular}{l}
\textbf{Виконали:} студенти 4 курсу, група МІ-41\\[4pt]
Петров Богдан\\
Мекшаков Артемій\\[10pt]
\textbf{Спеціальність:} Комп'ютерні науки\\
\end{tabular}
\end{flushright}
\vspace{1.5cm}
{Київ — #3\par}
\end{titlepage}
}


% SWOT-таблиця: #1 S, #2 W, #3 O, #4 T
\newcommand{\swot}[4]{%
\renewcommand{\arraystretch}{1.15}
\begin{tabularx}{\textwidth}{|Y|Y|}
\hline
\cellcolor{sgreen}\textbf{S — Сильні сторони} & \cellcolor{wred}\textbf{W — Слабкі сторони}\\
\hline
\cellcolor{sgreen!40}\begin{itemize}[leftmargin=1em,nosep]#1\end{itemize} &
\cellcolor{wred!40}\begin{itemize}[leftmargin=1em,nosep]#2\end{itemize}\\
\hline
\cellcolor{oblue}\textbf{O — Можливості} & \cellcolor{tyellow}\textbf{T — Загрози}\\
\hline
\cellcolor{oblue!40}\begin{itemize}[leftmargin=1em,nosep]#3\end{itemize} &
\cellcolor{tyellow!40}\begin{itemize}[leftmargin=1em,nosep]#4\end{itemize}\\
\hline
\end{tabularx}
}

\begin{document}
\maketitlepage{Звіт за етапом 1}{Вибір та обґрунтування теми командного проєкту. SWOT-аналіз варіантів}{2026}

\section{Мета етапу}
Визначити потенційні напрями розробки програмного продукту, проаналізувати їхню доцільність за допомогою SWOT-аналізу та обґрунтувати вибір теми командного проєкту.

\section{Критерії відбору тем}
Під час формування переліку варіантів команда керувалася такими обмеженнями:
\begin{itemize}[leftmargin=1.5em,nosep]
  \item реалістичний обсяг для команди з двох осіб протягом одного семестру (прототип до 18.11.2026);
  \item єдиний технологічний стек --- TypeScript, Next.js (клієнт), NestJS (сервер), PostgreSQL;
  \item наявність зрозумілої цільової аудиторії та реальної потреби, яку команда може перевірити на собі або на одногрупниках;
  \item можливість покрити продукт модульними та інтеграційними тестами, задокументувати API.
\end{itemize}

\section{Варіанти тем}

\subsection{Тема 1. «StudyFlow» --- планувальник навчальних дедлайнів для студентських груп}
\begin{tabularx}{\textwidth}{L{3.6cm} Y}
\toprule
\textbf{Призначення} & Вебзастосунок для ведення спільного календаря дедлайнів, лабораторних та командних задач у межах навчальної групи.\\
\midrule
\textbf{Цільова аудиторія} & Студенти університетів (насамперед групи та підгрупи, що виконують спільні проєкти), старости груп.\\
\midrule
\textbf{Проблема / потреба} & Дедлайни розпорошені між Classroom, месенджерами та усними домовленостями; студенти пропускають терміни, староста вручну дублює інформацію, відсутня єдина картина навантаження на тиждень.\\
\midrule
\textbf{Основні можливості} & Реєстрація та групи; створення предметів і дедлайнів (особистих та групових); канбан-дошка задач команди; календар тижня/місяця; нагадування на e-mail; відмітка виконання та прогрес групи; роль старости з правом масового додавання дедлайнів.\\
\bottomrule
\end{tabularx}

\vspace{6pt}
\swot{
  \item Команда є представником цільової аудиторії: вимоги перевіряються на собі.
  \item Предметна область проста: типовий CRUD, ролі, календар --- добре лягає на Next.js + NestJS.
  \item Легко декомпозувати на незалежні модулі (користувачі, групи, дедлайни, нагадування).
  \item Чітка реляційна модель даних, зручна для UML та тестування.
}{
  \item Низька технологічна новизна.
  \item Користувачам потрібно вручну вносити дедлайни; без інтеграцій цінність залежить від дисципліни групи.
  \item Нагадування потребують фонових задач (cron/черги) --- додаткова складність.
}{
  \item Інтеграція з Google Calendar / Classroom у майбутньому.
  \item Пілотне впровадження у власній групі МІ-41 та отримання реального зворотного зв'язку.
  \item Розширення до кількох факультетів, статистика навантаження для викладачів.
}{
  \item Конкуренція з Notion, Google Calendar, Trello.
  \item Ризик низького залучення: студенти можуть не підтримувати актуальність даних.
  \item Залежність від зовнішнього сервісу надсилання e-mail.
}

\subsection{Тема 2. «SplitMate» --- облік спільних витрат}
\begin{tabularx}{\textwidth}{L{3.6cm} Y}
\toprule
\textbf{Призначення} & Вебзастосунок для фіксації спільних витрат невеликих груп (гуртожиток, поїздка, спільна оренда) і автоматичного розрахунку, хто кому скільки винен.\\
\midrule
\textbf{Цільова аудиторія} & Студенти, що живуть разом або подорожують компаніями; молоді люди, які ділять побутові витрати.\\
\midrule
\textbf{Проблема / потреба} & Розрахунки «хто за що платив» ведуться у чатах і нотатках, губляться, спричиняють конфлікти; підрахунок балансу вручну трудомісткий.\\
\midrule
\textbf{Основні можливості} & Групи та учасники; додавання витрат із розподілом (порівну / частками / сумами); баланс кожного учасника; алгоритм мінімізації кількості переказів; історія та експорт у CSV; підтвердження погашення боргу.\\
\bottomrule
\end{tabularx}

\vspace{6pt}
\swot{
  \item Зрозуміла, добре формалізована бізнес-логіка; алгоритм спрощення боргів --- цікава задача для модульних тестів.
  \item Невелика кількість сутностей, швидкий MVP.
  \item Є реальна побутова потреба в команді та серед знайомих.
}{
  \item Робота з грошима вимагає точних обчислень (копійки, округлення) і підвищує ціну помилки.
  \item Обмежений простір для розвитку UI: продукт зводиться до кількох форм і таблиць.
  \item Без мобільного застосунку зручність нижча, ніж у конкурентів.
}{
  \item Підтримка валют, розпізнавання чеків, інтеграція з банківськими API.
  \item Монетизація через преміум-функції.
}{
  \item Сильні зрілі конкуренти (Splitwise, Tricount) з безкоштовними тарифами.
  \item Низька лояльність: користувачі переходять до звичних застосунків.
  \item Питання довіри до зберігання фінансових даних.
}

\subsection{Тема 3. «QuickPoll» --- миттєві опитування та голосування для занять}
\begin{tabularx}{\textwidth}{L{3.6cm} Y}
\toprule
\textbf{Призначення} & Сервіс для створення швидких опитувань, вікторин і голосувань із приєднанням за кодом або QR та відображенням результатів у реальному часі.\\
\midrule
\textbf{Цільова аудиторія} & Викладачі та студенти на лекціях і семінарах; організатори невеликих заходів.\\
\midrule
\textbf{Проблема / потреба} & Потрібен швидкий спосіб зібрати думку аудиторії або перевірити розуміння матеріалу без реєстрації учасників і складних інструментів.\\
\midrule
\textbf{Основні можливості} & Створення опитування з різними типами питань; доступ за коротким кодом без реєстрації; оновлення результатів у реальному часі (WebSocket); діаграми результатів; експорт відповідей; історія сесій для автора.\\
\bottomrule
\end{tabularx}

\vspace{6pt}
\swot{
  \item Ефектна демонстрація: результати оновлюються наживо.
  \item Низький поріг входу для користувачів (без реєстрації).
  \item Корисний для самих викладачів дисципліни.
}{
  \item Real-time (WebSocket) ускладнює архітектуру, тестування та деплой.
  \item Складніше забезпечити стабільність під час одночасного підключення багатьох учасників.
  \item Анонімний доступ ускладнює захист від накручування голосів.
}{
  \item Інтеграція з LMS університету, використання на конференціях.
  \item Розширення до гейміфікованих вікторин.
}{
  \item Дуже сильні конкуренти: Mentimeter, Kahoot, Slido, Google Forms.
  \item Ризик не встигнути стабілізувати real-time частину до дедлайну прототипу.
  \item Низька потреба: викладачі вже мають звичні інструменти.
}

\section{Порівняння варіантів}
Для об'єктивного вибору варіанти оцінено за п'ятьма критеріями за шкалою від 1 (погано) до 5 (відмінно). Ваги відображають пріоритети навчального проєкту: реалізованість у семестрі та відповідність вимогам курсу важливіші за ринковий потенціал.

\begin{table}[h]
\centering
\renewcommand{\arraystretch}{1.25}
\begin{tabularx}{\textwidth}{Y c c c c}
\toprule
\textbf{Критерій} & \textbf{Вага} & \textbf{StudyFlow} & \textbf{SplitMate} & \textbf{QuickPoll}\\
\midrule
Реалізованість командою з 2 осіб до 18.11 & 0{,}30 & 5 & 5 & 3\\
Відповідність вимогам курсу (UML, БД, тести, API-документація) & 0{,}25 & 5 & 4 & 3\\
Доступ до цільової аудиторії для перевірки гіпотез & 0{,}20 & 5 & 4 & 3\\
Потенціал розвитку після семестру & 0{,}15 & 4 & 3 & 4\\
Конкурентне середовище (5 --- слабке) & 0{,}10 & 3 & 2 & 1\\
\midrule
\textbf{Зважена оцінка} & 1{,}00 & \textbf{4{,}65} & \textbf{3{,}95} & \textbf{2{,}95}\\
\bottomrule
\end{tabularx}
\caption{Зважене порівняння варіантів тем}
\end{table}

\section{Обґрунтування вибору}
За результатами SWOT-аналізу та зваженого порівняння для подальшої розробки обрано \textbf{тему 1 --- «StudyFlow», планувальник навчальних дедлайнів для студентських груп}. Ключові аргументи:

\begin{enumerate}[leftmargin=1.5em]
  \item \textbf{Реалізованість.} Предметна область зводиться до добре відомих патернів (автентифікація, ролі, CRUD, календар), які повністю покриваються стеком Next.js + NestJS + PostgreSQL без ризикових технологій на зразок real-time. Обсяг MVP реалістичний для двох розробників за чотири тижневі спринти.
  \item \textbf{Відповідність вимогам курсу.} Реляційна модель (користувач, група, предмет, дедлайн, задача, нагадування) дає змістовну схему БД та UML-діаграми; бізнес-логіка (права ролей, розрахунок прогресу, планування нагадувань) природно покривається 20+ модульними тестами; REST API документується через Swagger.
  \item \textbf{Доступ до аудиторії.} Команда та її одногрупники є безпосередніми користувачами, тому гіпотези можна перевірити на реальній групі ще в межах семестру, а зворотний зв'язок використати в рефлексіях спринтів.
  \item \textbf{Рівномірний розподіл роботи.} Продукт має однаково вагомі серверну (API, ролі, нагадування, БД) та клієнтську (календар, канбан, адаптивний інтерфейс) частини, що відповідає розподілу ролей у команді.
  \item \textbf{Керованість загроз.} Загрозу конкуренції з універсальними інструментами знижує вузька спеціалізація на навчальному процесі (предмети, групи, роль старости); залежність від e-mail-сервісу ізолюється окремим модулем нагадувань.
\end{enumerate}

Тема 2 («SplitMate») визнана прийнятною резервною альтернативою; тема 3 («QuickPoll») відхилена через підвищений технічний ризик real-time складової та насичений ринок.

\section{Висновки}
Запропоновано три варіанти теми програмного продукту, для кожного складено короткий опис і SWOT-аналіз, проведено зважене порівняння. Обґрунтовано вибір теми «StudyFlow». На наступному етапі (ініціація проєкту, до 27.10.2026) команда підготує Lean Canvas, карту продукту та заведе дошку проєкту на платформі управління.

\end{document}
